\section{Introduction}

DNA methylation is a covalent epigenetic modification that contributes to transcriptional
regulation, chromatin organization, genome stability, and the maintenance of cellular identity.
At CpG dinucleotides, methylation is commonly quantified as a continuous beta value using
array-based assays that profile hundreds of thousands of loci \citep{Bibikova2011,Pidsley2016}.
The biological meaning of a methylation measurement depends on both genomic and cellular
context: regulatory CpGs can reflect altered gene activity, whereas tissue-specific methylation
programs capture differentiation and cell composition \citep{Jones2012,Roadmap2015}. A
methylation profile therefore represents an integrated molecular state rather than a collection
of independent biomarkers.

The stability of DNA methylation, its accessibility across biological specimens, and its
responsiveness to development, ageing, environmental exposure, and disease have made it a
valuable substrate for biomarker discovery. Epigenetic clocks illustrate this potential. Early
regularized models estimated chronological age across tissues \citep{Horvath2013}, while later
clocks were developed to capture mortality risk, health-related ageing, and the pace of
physiological change \citep{Levine2018,Lu2019,Belsky2022,HigginsChen2022}. These models remain
strong and interpretable predictors, but each is optimized for a predefined endpoint; biological
structure shared across methylation datasets must therefore be learned repeatedly for different
phenotypes.

This task-specific paradigm is poorly matched to the scale and heterogeneity of public
methylation data. Existing profiles span tissues, cohorts, disease states, and array platforms,
while high-quality phenotype labels are available for only a subset of samples and differ among
studies. Moreover, platform-defined CpG panels overlap only partially, so models must often
operate on different subsets of the methylome. Conventional linear predictors perform well when
the endpoint and CpG panel are fixed, but their predominantly additive formulation does not
directly model context-dependent relationships among loci. Deep supervised models can capture
such relationships, yet training a separate model for each application leaves the much larger
unlabelled corpus unused and can generalize poorly from small labelled cohorts. These constraints
motivate a foundation-model strategy that separates large-scale representation learning from
task-specific supervision.

Self-supervised pretraining has enabled reusable representations in language, genomics, and
single-cell transcriptomics \citep{Devlin2019,Ji2021,Yang2022,Cui2024}. In DNA methylation,
MethylGPT established that a Transformer can be pretrained on a large collection of human
methylomes and adapted to biological prediction tasks \citep{Ying2024}. MethylGPT combines
masked CpG prediction with whole-profile reconstruction from a global CLS representation,
thereby encouraging that representation to retain sample-level information. Whole-profile
decoding, however, does not explicitly require the representation of a sample to remain stable
when different CpG subsets are observed, nor does it isolate recovery of loci that were entirely
absent from the encoder input. These properties are important when a pretrained representation
must transfer across incomplete profiles and non-identical measurement panels.

Here we present \method{}, a self-supervised foundation model for human DNA methylation.
\method{} was trained from scratch, without phenotype labels, on 169{,}120 profiles represented
at 49{,}156 curated CpG sites, using the large public methylation resource and CpG panel assembled
for MethylGPT \citep{Ying2024}. Each token combines CpG identity with its quantitative
methylation value. A Transformer encoder integrates these measurements, while genomic-rank
rotary position encoding preserves chromosomal order and a dedicated CLS token provides a
global representation for downstream adaptation \citep{Su2024}.

To learn this representation, we adapt the whole cell expression decoder objective introduced in
BMFM-RNA \citep{Dandala2025} to methylation as the Whole-CpG Encoder--Decoder (WCED). Two
independently sampled partial views are generated from each methylation profile and processed by
a shared encoder. For each view, a decoder predicts the measured profile from the CLS
bottleneck, but the reconstruction loss is evaluated only at valid CpGs excluded from that
encoder input, so the objective cannot be satisfied by copying the input forward. A
complementary response to the same problem is offered by scConcept, which discards feature
reconstruction altogether and pretrains cell representations purely by contrasting disjoint
gene-panel views \citep{Bahrami2025}. We adopt that diagnosis but not that remedy: rather than
replacing reconstruction, we retain it in bottleneck form and add an InfoNCE objective that
aligns CLS representations of the two views from the same profile while distinguishing
representations from different profiles \citep{Oord2018}. WCED therefore applies two distinct
pressures to the same embedding, which must both support recovery of unobserved methylation
values and remain stable across CpG subsets.

We show that independently sampled CpG views preserve sample-specific identity and that the CLS
bottleneck supports accurate reconstruction of CpGs withheld from the encoder. Without phenotype
supervision, the pretrained representation captures substantial age- and tissue-associated
variation, and fine-tuning for chronological-age prediction outperforms both a matched MethylGPT
pipeline and an ElasticNet baseline on a common held-out cohort. At the same time, frozen
representations encode strong study-specific structure, and transfer of age information to
entirely unseen studies is heterogeneous. Together, these results establish \method{} as a
reusable DNA methylation foundation model while identifying cohort heterogeneity as a central
challenge for cross-study generalization.

% Internal drafting note (not printed): add additional completed downstream tasks before submission.
